\documentclass[12pt]{article}
\usepackage[utf8]{inputenc}
\usepackage[T1]{fontenc}
\usepackage{hyperref}
\title{Tat, Muslim Language Resources for K-12 Curriculum Development}
\author{Compiled by Claude Sonnet 5 (claude-sonnet-5) for Andr\'as Kornai}
\date{}
\begin{document}
\maketitle
\section{Introduction}
Muslim Tat (ISO 639-3 ttt; Glottocode musl1236) is an Iranian (Indo-European) language of Azerbaijan and Dagestan, Russia. Speaker estimates vary widely: about 22,800 in Azerbaijan (2011) and about 2,000 in Dagestan (2010) in one source, and 34,000 in another. The language is shifting toward Azerbaijani. Search found background on status and writing but no current Muslim Tat textbook or modern dictionary. The Judeo-Tat (Juhuri) materials that appear in results are for a different, related language and should not be confused with it.
\section{Primary Educational Resources}
No textbook was found. Sources report that Tat classes have been started in several schools in the Quba region, using an alphabet based on the Azerbaijani Latin alphabet.
\section{Linguistic and Reference Materials}
\begin{itemize}
\item Wikipedia, Tat language (Caucasus): \url{https://en.wikipedia.org/wiki/Tat_language_%28Caucasus%29}
\item Wikipedia, Tat alphabet: \url{https://en.wikipedia.org/wiki/Tat_alphabet}
\item Omniglot, Tat: \url{https://www.omniglot.com/writing/tat.htm}
\item Red Book of the Peoples of the Russian Empire, Tats: \url{https://www.eki.ee/books/redbook/tats.shtml}
\end{itemize}
Vladimir Miller's \emph{Tati Studies} (1905) is named in search summaries as containing a Tat--Russian dictionary. I did not see a link to the text.
\section{Religious and Community Resources}
None found in search.
\section{Digital and Audio Resources}
None found in search. Sources mention a small online revival movement but give no specific resources.
\section{Curriculum Development Context}
Tat is taught as a subject in some schools (Quba region, Azerbaijan), and it has official status in Dagestan, where the script is Cyrillic. Azerbaijan has funded minority-language development since 1996. Curriculum work would have to reconcile the Latin-based and Cyrillic-based writing systems and would probably need new materials.
\end{document}